\section*{अध्याय \ref{ch:b2:diatomic-mos} --- द्विपरमाणुक अणुओं के आण्विक कक्षक}

\begin{solution}{exo:b2:diatomic-mos:1}
$E_+ = (-13.6 - 9.0)/1.60 = \qty{-14.1}{eV}$; $E_- = (-13.6 + 9.0)/0.40 = \qty{-11.5}{eV}$।
$\psi_+$ में दो इलेक्ट्रॉन: $-28.25$, जबकि $2 \times (-13.6) = \qty{-27.2}{eV}$; इस स्थूल मॉडल में
\qty{1.05}{eV} से बद्ध।
\end{solution}

\begin{solution}{exo:b2:diatomic-mos:2}
\ce{He2}: $(2 - 2)/2 = 0$; \ce{He2+}: $(2 - 1)/2 = 0.5$; \ce{Li2}: 1; \ce{B2}: छह संयोजकता
इलेक्ट्रॉन, $\sigma_{2s}^2\sigma_{2s}^{*2}\pi_{2p}^2$, \omterm{def:b2:diatomic-mos:bond-order}{आबंध कोटि} 1।
\end{solution}

\begin{solution}{exo:b2:diatomic-mos:3}
\ce{B2} (दो $\pi_{2p}$ कक्षकों में दो अयुग्मित इलेक्ट्रॉन, s और p के मिश्रण के साथ) और
\ce{O2} ($\pi^*$ में दो)। \ce{C2}, \ce{N2}, \ce{F2} \omterm{def:b2:diatomic-mos:magnetism}{प्रतिचुंबकीय} हैं।
\end{solution}

\begin{solution}{exo:b2:diatomic-mos:4}
शून्येतर: (1s, $2\mathrm p_z$), ($2\mathrm p_x$, $2\mathrm p_x$), (2s, $2\mathrm p_z$)।
सममिति से शून्य: (1s, $2\mathrm p_x$) और ($2\mathrm p_x$, $2\mathrm p_y$)।
\end{solution}

\begin{solution}{exo:b2:diatomic-mos:5}
आयनन एक आबंधी इलेक्ट्रॉन हटाता है: \omterm{def:b2:diatomic-mos:bond-order}{आबंध कोटि} 3 के स्थान पर 2.5, अतः लंबा और
दुर्बल आबंध, जैसा $D_0(\ce{N2+}) = \qty{8.71}{eV} < D_0(\ce{N2}) = \qty{9.76}{eV}$ दिखाता है।
\end{solution}

\begin{solution}{exo:b2:diatomic-mos:6}
ऑक्सीजन के कक्षक नीचे हैं (ऑक्सीजन अधिक विद्युत-ऋणात्मक है):
\cref{thm:b2:diatomic-mos:heteronuclear} से \omterm{def:b2:diatomic-mos:bonding}{आबंधी कक्षक}
निचले परमाणु, ऑक्सीजन, पर भारित होते हैं, और ऊपर के कक्षक कार्बन पर। उच्चतम भरा
कक्षक, ऑक्सीजन से दूर इंगित करता अधिकांशतः कार्बन-लक्षण वाला एक $\sigma$ कक्षक,
कार्बन पर एकाकी युग्म है।
\end{solution}

\begin{solution}{exo:b2:diatomic-mos:7}
केवल फ़्लुओरीन के $2\mathrm p_z$ (अक्ष के अनुदिश) की सममिति हाइड्रोजन के 1s जैसी है:
वे F पर भारित एक $\sigma$ कक्षक और H पर भारित एक $\sigma^*$ बनाते हैं। फ़्लुओरीन के
$2\mathrm p_x$, $2\mathrm p_y$ (शून्य अतिव्यापन) और उसका नीचा 2s (ऊर्जा में बहुत दूर)
अनाबंधी रहते हैं: तीन एकाकी युग्म। आठ संयोजकता इलेक्ट्रॉन 2s, $\sigma$
और दो अनाबंधी p को भरते हैं।
\end{solution}

\begin{solution}{exo:b2:diatomic-mos:8}
$\bar\alpha = \qty{-12}{eV}$, $\Delta = \qty{4}{eV}$, $\sqrt{4 + 9} = \qty{3.61}{eV}$:
$E = \qty{-15.6}{eV}$ और \qty{-8.4}{eV}। निचले स्तर के लिए $(\alpha_B - E)c_B = -\beta
c_A$: $1.61c_B = 3.0c_A$, $c_A/c_B = 0.535$, $c_B^2 = 1/(1 + 0.286) = 0.78$।
\end{solution}

\begin{solution}{exo:b2:diatomic-mos:9}
आठ संयोजकता इलेक्ट्रॉन: $\sigma_{2s}^2\sigma_{2s}^{*2}\pi_{2p}^4$, जबकि $\sigma_{2p}$ कक्षक
(मिश्रण द्वारा $\pi_{2p}$ से ऊपर धकेला गया) खाली है। \omterm{def:b2:diatomic-mos:bond-order}{आबंध कोटि} $(6 - 2)/2 = 2$, दो
भरे $\pi$ कक्षकों से।
\end{solution}

\begin{solution}{exo:b2:diatomic-mos:10}
NO, ग्यारह संयोजकता इलेक्ट्रॉन, एक $\pi^*$ में: \omterm{def:b2:diatomic-mos:bond-order}{आबंध कोटि} 2.5। \ce{NO+}, दस: \omterm{def:b2:diatomic-mos:bond-order}{आबंध कोटि}
3। $D_0(\ce{NO+}) = 6.51 + 13.62 - 9.26 = \qty{10.87}{eV}$: NO से कहीं प्रबल, क्योंकि
हटाया गया इलेक्ट्रॉन प्रतिआबंधी था।
\end{solution}

\begin{solution}{exo:b2:diatomic-mos:11}
$E_+ + E_- = \dfrac{(\alpha + \beta)(1 - S) + (\alpha - \beta)(1 + S)}{1 - S^2} = \dfrac{2(\alpha -
\beta S)}{1 - S^2}$। यह $2\alpha$ से अधिक है जब $\alpha - \beta S > \alpha - \alpha S^2$, अर्थात्
$\beta < \alpha S$ ($-S < 0$ से भाग देकर)। चार इलेक्ट्रॉनों के साथ, हर कक्षक में दो,
कुल ऊर्जा अलग कक्षकों से ऊपर है: दो भरे कक्षक एक-दूसरे को प्रतिकर्षित करते हैं।
\end{solution}

\begin{solution}{exo:b2:diatomic-mos:12}
\omterm{def:b2:diatomic-mos:bond-order}{आबंध कोटियाँ} 2.5, 2, 1.5, 1: इसी क्रम में आबंध लंबे और दुर्बल होते हैं। हाइड्रोजन
परॉक्साइड में परॉक्साइड O--O एकल आबंध (\ce{O2^2-} इकाई) है; \ce{KO2} में
सुपरऑक्साइड आयन \ce{O2-} है, जो \omterm{def:b2:diatomic-mos:magnetism}{अनुचुंबकीय} है।
\end{solution}

\begin{solution}{pb:b2:diatomic-mos:1}
\textbf{1.} बारह।
\textbf{2.} $\sigma_{2s}^2\,\sigma_{2s}^{*2}\,\sigma_{2p}^2\,\pi_{2p}^4\,\pi_{2p}^{*2}$, दोनों $\pi^*$
इलेक्ट्रॉन भिन्न कक्षकों में।
\textbf{3.} $(8 - 4)/2 = 2$।
\textbf{4.} दो: \ce{O2} \omterm{def:b2:diatomic-mos:magnetism}{अनुचुंबकीय} है।
\textbf{5.} वह सभी इलेक्ट्रॉनों को युग्मित करती है, जबकि उनमें से दो दो अपभ्रष्ट
$\pi^*$ कक्षकों में अकेले रहते हैं।
\textbf{6.} $D_0 = 2 \times 246.79 = \qty{493.6}{kJ/mol}$, \qty{5.12}{eV}।
\textbf{7.} \ce{O2+}: $\pi^{*1}$, \omterm{def:b2:diatomic-mos:bond-order}{आबंध कोटि} 2.5।
\textbf{8.} \ce{O2-}: $\pi^{*3}$, \omterm{def:b2:diatomic-mos:bond-order}{आबंध कोटि} 1.5।
\textbf{9.} \ce{O2^2-}: $\pi^{*4}$, \omterm{def:b2:diatomic-mos:bond-order}{आबंध कोटि} 1।
\textbf{10.} \ce{O2+} एक, \ce{O2-} एक, \ce{O2^2-} कोई नहीं।
\textbf{11.} \ce{O2+} $<$ \ce{O2} $<$ \ce{O2-} $<$ \ce{O2^2-}।
\textbf{12.} परॉक्साइड आयन।
\textbf{13.} \ce{O2+}, \ce{O2} और \ce{O2-}।
\textbf{14.} एक $\pi^*$ कक्षक से।
\textbf{15.} प्रतिआबंधी $\pi^*$ कक्षक परमाणु के 2p स्तर से ऊपर है: उसका
इलेक्ट्रॉन कम बद्ध है।
\textbf{16.} \ce{O2} से $\ce{O} + \ce{O+} + \mathrm e^-$ तक दो पथ: वियोजन फिर
एक परमाणु का आयनन, $D_0(\ce{O2}) + I(\ce{O})$; या अणु का आयनन फिर आयन का वियोजन,
$I(\ce{O2}) + D_0(\ce{O2+})$। बराबर: $D_0(\ce{O2+}) = 5.12 + 13.62 - 12.07 = \qty{6.66}{eV}$।
\textbf{17.} $D_0(\ce{O2}) = \qty{5.12}{eV}$ से बड़ा: प्रतिआबंधी इलेक्ट्रॉन हटाने से
आबंध प्रबल होता है (कोटि 2.5)।
\textbf{18.} $D_0(\ce{N2}) = 2 \times 470.82/96.485 = \qty{9.76}{eV}$; $D_0(\ce{N2+}) = 9.76 +
14.53 - 15.58 = \qty{8.71}{eV}$।
\textbf{19.} \ce{N2} का उच्चतम भरा कक्षक आबंधी है, परमाणु के 2p स्तर से
नीचे।
\textbf{20.} जब हटाया गया इलेक्ट्रॉन किसी \omterm{def:b2:diatomic-mos:bonding}{प्रतिआबंधी कक्षक} से आता है।
\textbf{21.} हुंड का नियम (अपभ्रष्ट कक्षकों में समांतर चक्रणों की अधिकतम संख्या)।
\textbf{22.} दोनों $\pi^*$ इलेक्ट्रॉन एक ही कक्षक में, चक्रण युग्मित।
\textbf{23.} एक ही कक्षक में दो इलेक्ट्रॉन एक-दूसरे को अधिक प्रतिकर्षित करते हैं, और समांतर
चक्रणों का विनिमय स्थायीकरण खो देते हैं।
\textbf{24.} \omterm{def:b2:diatomic-mos:bond-order}{आबंध कोटि} अब भी 2; सभी इलेक्ट्रॉन युग्मित: \omterm{def:b2:diatomic-mos:magnetism}{प्रतिचुंबकीय}।
\textbf{25.} \ce{O2+} की \omterm{def:b2:diatomic-mos:bond-order}{आबंध कोटि} $\boldsymbol{2.5}$ और वियोजन ऊर्जा
$\boldsymbol{\approx \qty{6.66}{eV}}$ है, जबकि \ce{O2} के लिए 2 और \qty{5.12}{eV}।
\end{solution}
